\documentclass[11pt]{amsart}

\usepackage[margin=1in]{geometry}
\usepackage{amsmath,amssymb,amsthm}
\usepackage{booktabs}
\usepackage{longtable}
\usepackage{array}
\usepackage{siunitx}
\usepackage{hyperref}
\usepackage{xcolor}
\usepackage{listings}

\hypersetup{
  colorlinks=true,
  linkcolor=blue!50!black,
  citecolor=blue!50!black,
  urlcolor=blue!50!black
}

\newtheorem{lemma}{Lemma}
\newtheorem{theorem}{Theorem}

\newcommand{\pcount}{\pi}
\newcommand{\li}{\operatorname{li}}

\title[Ramanujan Primes Beyond $10^{19}$]{%
Ramanujan Primes Beyond $10^{19}$:\\
A Certified Extension of OEIS A181671 via a 128-bit\\
Segmented Sieve and Bracketed Analytic Tail Bounds}

\author{Gaurav Vyas}
\thanks{Email: \texttt{gaurav.vyas.1729@gmail.com}}
\email{gaurav.vyas.1729@gmail.com}

\date{\today}

\begin{document}

\begin{abstract}
A Ramanujan prime is the smallest integer $R_n$ such that
$\pcount(x) - \pcount(x/2) \geq n$ for all $x \geq R_n$, where $\pcount$ is
the ordinary prime-counting function. The count of Ramanujan primes below
$10^k$, tabulated as OEIS sequence A181671, has published values only
through $k=17$; two further terms ($k=18,19$) exist as unpublished prior
computations. We present a certified extension of this sequence to
$k=20,21,22$, and (in progress at time of writing) $k=23$, computed on
consumer hardware with no institutional resources. The method combines
a bracketing lemma that reduces the certification of $f(x) = \pcount(x) -
\pcount(x/2)$ on an entire interval to exact evaluations of $\pcount$ at
finitely many endpoints, a custom 128-bit segmented sieve for the region
immediately above $10^{k}$ (since general-purpose sieve libraries are
limited to $2^{64}$), and the explicit analytic bounds of
Dusart~\cite{dusart2010} and Johnston~\cite{johnston2022} to certify the
tail beyond the sieved and bracketed region without further computation.
Every new term's defining anchors $\pcount(Q)$ and $\pcount(Q/2)$ are
cross-verified by two independent algorithms (Gourdon and
Deleglise--Rivat) implemented in the \texttt{primecount}
library~\cite{primecount}, and selected results are further verified
against an independent published reference for $\pcount(x)$ and by direct
Miller--Rabin primality testing that bypasses the custom sieve entirely.
We report the full implementation, the RAM- and CPU-aware scheduling
required to complete these computations reliably on hardware with
16--32\,GB of RAM, and an empirical time-complexity analysis showing the
method scales as $\sim\!3.5$--$3.7\times$ per decade of $Q$ once
measurement artifacts from interrupted runs are removed from the data
--- better than the na\"ive $Q^{2/3}$ asymptotic predicted for the
underlying prime-counting algorithm.
\end{abstract}

\maketitle

\section{Introduction}

Ramanujan's 1919 paper~\cite{ramanujan1919} proved a strengthening of
Bertrand's postulate: for every $n \geq 1$ there exists a least integer
$R_n$ such that
\[
  \pcount(x) - \pcount(x/2) \geq n \qquad \text{for all } x \geq R_n.
\]
These $R_n$ are now called Ramanujan primes. The number of Ramanujan
primes not exceeding $10^k$, denoted $a(k)$, is tabulated in the OEIS as
sequence A181671~\cite{oeisA181671}. As of this writing the OEIS entry
publishes verified values only for $k \leq 17$; the values $a(18)$ and
$a(19)$ circulate as unpublished, self-computed extensions with no
independent verification on record.

The obstacle to extending this sequence further is entirely
computational, not mathematical: the natural method is to sieve the
interval $[10^k, 10^k + W]$ for some window $W$ large enough to certify
that the running count of $f(x) = \pcount(x) - \pcount(x/2)$ never dips
below $a(k)-1$ again. Under the na\"ive approach $W$ grows with $\sqrt{Q}$,
and general-purpose prime-sieving libraries such as
\texttt{primesieve}~\cite{primesieve} are limited to 64-bit integers, so
this approach becomes infeasible well before $Q = 10^{20}$: at that scale
the certification window alone spans on the order of $10^{12}$
integers, entirely above the $2^{64}$ boundary.

This paper reports a certified extension of $a(k)$ to $k=20,21,22$, with
$k=23$ in progress at the time of writing (a placeholder is left in
Table~\ref{tab:results} pending completion; see Section~\ref{sec:a23}).
Every term is accompanied by a machine-checkable certificate and passes
multiple layers of independent verification described in
Section~\ref{sec:verification}. The entire computation was carried out
on two personal machines (Section~\ref{sec:hardware}) with no cluster,
cloud, or institutional compute access.

\section{Background}
\label{sec:background}

\subsection{Ramanujan primes and the bracketing identity}

Write $f(x) = \pcount(x) - \pcount(x/2)$. Since $\pcount$ is
nondecreasing, $f$ is nondecreasing wherever it is evaluated at integer
multiples that avoid parity artifacts, and in particular:

\begin{lemma}[Bracketing lemma]
\label{lem:bracket}
For $a \leq y \leq b$,
\[
  f(y) \;\geq\; \pcount(a) - \pcount(b/2) \;=\; f(a) - \bigl[\pcount(b/2) -
  \pcount(a/2)\bigr].
\]
\end{lemma}

\begin{proof}
$\pcount$ is nondecreasing, so $\pcount(y) \geq \pcount(a)$ and
$\pcount(y/2) \leq \pcount(b/2)$ for $a \leq y \leq b$; subtract.
\end{proof}

This is elementary, but it is the entire mechanism that makes evaluation
at $Q = 10^{20}$ and beyond tractable: a single exact evaluation of
$\pcount$ at each of two endpoints lower-bounds $f$ across the whole
interval $[a,b]$, with no sieving of the interior required. Because $f$
drifts upward at rate $\sim 1/(2\log Q)$ in practice, the bound survives
\emph{geometric} growth of the interval width, so only
$O(\log(W/D))$ grid points are needed to bracket a window of width $W$
starting from an initial exactly-sieved region of width $D$.

\subsection{Analytic tail bounds}

Grid bracketing alone still requires evaluating $\pcount$ out to some
finite $Y$ before the interval can be declared safe. Beyond $Y$, we use
two published unconditional (RH-independent) explicit bounds on
$\pcount$:

\begin{itemize}
  \item \textbf{Johnston (2022)}~\cite{johnston2022}: using a
  computer-verified region of the Riemann hypothesis up to height
  $3\times10^{12}$,
  \[
    |\pcount(x) - \li(x)| < \frac{\sqrt{x}\,\log x}{8\pi}
    \qquad \text{for } 2657 \leq x \leq 1.101\times10^{26}.
  \]
  Applying this at both $x=y$ and $x=y/2$ (requiring $y \geq 5314$) gives
  the certifiable lower bound used in this work,
  \[
    G(y) \;=\; \li(y) - \li(y/2) -
    \frac{\sqrt{y}\log y + \sqrt{y/2}\log(y/2)}{8\pi} \;\leq\; f(y).
  \]
  \item \textbf{Dusart (2010)}~\cite{dusart2010}, Theorem 6.9: for
  $x \geq 88{,}783$,
  \[
    \pcount(x) \geq \frac{x}{\ln x}\left(1+\frac1{\ln x}+\frac2{\ln^2 x}\right),
  \]
  and for $x \geq 2{,}953{,}652{,}287$,
  \[
    \pcount(x) \leq \frac{x}{\ln x}\left(1+\frac1{\ln x}+\frac{2.334}{\ln^2 x}\right).
  \]
  These cover $y \geq 1.101\times10^{26}$, beyond the Johnston bound's
  region of validity, and are what ultimately bound this method's reach:
  since the Johnston bound requires $y \leq 1.101\times10^{26}$, and the
  method needs it to hold at $y=Q$ for the largest term certified, the
  sequence becomes uncertifiable by this pipeline once $Q$ approaches
  that limit --- which occurs at $a(26)$ (Section~\ref{sec:ceiling}).
\end{itemize}

Both bounds were transcribed into the implementation from their
respective published papers and re-verified against the primary sources
as part of this work (Section~\ref{sec:verification}).

\section{Implementation}
\label{sec:implementation}

\subsection{Why a custom 128-bit sieve was necessary}

The exact-count region $[Q, Q+D]$ immediately above $Q$ must be sieved
directly rather than bracketed, since the bracketing lemma alone gives no
information about where inside an interval $f$ might dip. At
$Q=10^{20}$ this region already lies entirely above $2^{64} \approx
1.8\times10^{19}$, outside the range of \texttt{primesieve} and every
off-the-shelf sieve library available at the time of writing. We wrote a
custom 128-bit segmented sieve, \texttt{sieve128.cpp}, implementing:

\begin{itemize}
  \item a standard segmented sieve of Eratosthenes over 128-bit unsigned
  integers, using \texttt{primesieve} internally only to generate the
  base primes up to $\sqrt{Q+D}$ (which remain well within 64-bit range
  even at $Q=10^{23}$, since $\sqrt{10^{23}} \approx 3.16\times10^{11}$);
  \item an exact walk of $f$ across the segmented region, tracking the
  running minimum and the offset at which it occurs, so that the
  certificate records not just the final count but the exact point (if
  any) where $f$ dips within the sieved window;
  \item cross-validation against \texttt{primesieve} below $2^{64}$
  (exact agreement on prime lists, including across the $2^{64}$
  boundary) and against \texttt{sympy}'s primality routines above
  $2^{64}$, both performed as part of this project's own validation
  suite before any new term was attempted.
\end{itemize}

\subsection{Software architecture}

The pipeline consists of three components:

\begin{enumerate}
  \item \texttt{sieve128.cpp/.exe} --- the 128-bit segmented sieve and
  exact walk, described above.
  \item \texttt{primecount} (upstream, unmodified except for build
  configuration) --- provides exact, arbitrary-precision $\pcount(x)$ via
  two independently implemented algorithms, Gourdon's method and the
  Deleglise--Rivat method, used here purely as a cross-check against each
  other rather than for speed.
  \item \texttt{ramanujan128.py} --- the orchestration driver: computes
  and caches every $\pcount$ value obtained, invokes the sieve walk near
  each new $Q$, evaluates the analytic tail bounds, and assembles a
  JSON certificate per term recording every anchor, checkpoint, and
  verification outcome.
\end{enumerate}

\subsection{Crash-safe persistent caching}

Because a single anchor evaluation at $Q=10^{23}$ can run for several
hours (Section~\ref{sec:performance}), and because this project's
compute budget consisted of ordinary desktop/laptop machines subject to
reboots, power interruptions, and resource contention with unrelated
workloads, every exact $\pcount$ value ever computed is persisted to
\texttt{pi\_cache.json} using an atomic write pattern: results are
written to a temporary file and then renamed over the target, so that a
process kill or power loss mid-write can never leave the cache in a
corrupted state. This was validated directly by simulating fifteen
mid-write process kills, all of which left the cache file intact. In
practice, this pipeline survived at least one full unattended power
interruption and two accidental process terminations during the
computation of the terms reported here, in every case resuming from
exactly the last completed value with no lost work beyond whatever
single call was in flight at the moment of interruption.

\subsection{RAM- and CPU-aware adaptive parallelism}

Measured memory footprint per concurrent \texttt{primecount} call grows
substantially faster than the algorithm's own $\sqrt[3]{x}$ working-set
prediction: on the hardware described in Section~\ref{sec:hardware},
peak RAM per call was measured directly at approximately 0.36\,GB at
$Q=5\times10^{20}$, 1.26\,GB at $Q=10^{22}$, and 5.9\,GB at $Q=10^{23}$
--- an empirical scaling exponent of roughly $0.5$--$0.7$ in $Q$, not the
$1/3$ the underlying algorithm's complexity would suggest. Running
multiple checkpoint evaluations concurrently is necessary to make full
use of available CPU cores (a single \texttt{primecount} call was
measured to use only $\sim\!4$ of 12 available logical cores, since its
dominant phases are not fully parallel), but naively launching as many
concurrent calls as there are cores risks exhausting system RAM and
crashing the host --- which was observed directly once during this
project, when RAM contention with an unrelated concurrent workload
coincided with a full system crash. The driver therefore queries live
available RAM before sizing each batch of concurrent evaluations,
computes an expected per-call footprint from the fitted empirical curve
above with a safety margin, and caps concurrency accordingly, trading
some throughput for reliability on hardware without dedicated headroom.

\subsection{Cross-machine reproduction and migration}
\label{sec:hardware}

Terms $a(1)$ through $a(22)$ were computed on a laptop-class machine (13th
Gen Intel Core i5-13420H, 8 cores/12 threads, 16\,GB RAM, NVIDIA GeForce
RTX 4050 Laptop GPU, not used for this computation). Term $a(23)$ is being
computed on a separate desktop machine (12th Gen Intel Core i5-12400F, 6
cores/12 threads, 32\,GB RAM, NVIDIA GeForce RTX 3060, likewise not used
for this computation), after the laptop's available RAM proved
insufficient for $Q=10^{23}$: a single \texttt{primecount} call at that
scale was measured to require approximately 5.9\,GB, leaving under 1\,GB
of headroom on the 13.7\,GB effectively available on the laptop for the
many hours the computation would require, and the attempt was
deliberately cancelled rather than risked. The persistent JSON cache and
all source files transferred verbatim between machines (verified by
SHA-256 hash equality before resuming), so no computation already
completed was repeated.

As a byproduct, this migration produced a genuine independent
cross-machine reproduction: $a(18)$ and $a(19)$ were recomputed from the
persisted cache on the desktop machine using a from-scratch driver
invocation, and matched the laptop-computed values exactly on every
intermediate quantity (Table~\ref{tab:crossmachine}).

\begin{table}[h]
\centering
\caption{Cross-machine independent reproduction of $a(18)$, $a(19)$.}
\label{tab:crossmachine}
\begin{tabular}{lrr}
\toprule
Quantity & $k=18$ & $k=19$ \\
\midrule
$\pcount(10^k)$ & 24{,}739{,}954{,}287{,}740{,}860 & 234{,}057{,}667{,}276{,}344{,}607 \\
$\pcount(10^k/2)$ & 12{,}585{,}956{,}566{,}571{,}620 & 118{,}959{,}989{,}688{,}273{,}472 \\
$a(k)$ & 12{,}153{,}997{,}721{,}169{,}239 & 115{,}097{,}677{,}588{,}071{,}134 \\
\bottomrule
\end{tabular}
\end{table}

\section{Results}
\label{sec:results}

{\footnotesize
\begin{longtable}{@{} r r r r c @{}}
\caption{$a(k)$ for $k=1,\dots,23$. Times for $k \leq 18$ and $k=21,22$
are clean, single-pass, cache-bypassed measurements; $k=19,20$ are
genuine first-time measurements with no interruption. See
Section~\ref{sec:performance} for why some originally recorded times
were superseded by fresh recomputation. Status codes are defined below
the table.}
\label{tab:results} \\
\toprule
$k$ & $a(k)$ & time & chk. & status \\
\midrule
\endfirsthead
\toprule
$k$ & $a(k)$ & time & chk. & status \\
\midrule
\endhead
1  & 1 & 0.64\,s & 0 & A \\
2  & 10 & 0.38\,s & 0 & A \\
3  & 72 & 0.38\,s & 0 & A \\
4  & 559 & 0.38\,s & 0 & A \\
5  & 4{,}459 & 0.38\,s & 0 & A \\
6  & 36{,}960 & 0.37\,s & 0 & A \\
7  & 316{,}066 & 0.36\,s & 0 & A \\
8  & 2{,}760{,}321 & 0.36\,s & 0 & A \\
9  & 24{,}491{,}666 & 0.35\,s & 0 & A \\
10 & 220{,}098{,}288 & 0.34\,s & 0 & A \\
11 & 1{,}998{,}400{,}235 & 0.35\,s & 0 & A \\
12 & 18{,}299{,}775{,}876 & 0.41\,s & 1 & A \\
13 & 168{,}773{,}875{,}190 & 0.59\,s & 3 & A \\
14 & 1{,}566{,}017{,}986{,}235 & 1.40\,s & 6 & A \\
15 & 14{,}606{,}736{,}768{,}049 & 4.41\,s & 8 & A \\
16 & 136{,}860{,}923{,}837{,}558 & 16.57\,s & 10 & A \\
17 & 1{,}287{,}462{,}389{,}890{,}262 & 91.81\,s & 13 & A \\
18 & 12{,}153{,}997{,}721{,}169{,}239 & 5\,m\,43\,s & 15 & B \\
19 & 115{,}097{,}677{,}588{,}071{,}134 & 11\,m\,18\,s & 17 & B \\
20 & 1{,}093{,}039{,}678{,}770{,}734{,}297 & 63\,m\,59\,s & 16 & C \\
21 & 10{,}406{,}559{,}368{,}229{,}726{,}028 & 4\,h\,31\,m\,47\,s & 18 & D \\
22 & 99{,}306{,}360{,}875{,}818{,}676{,}888 & 14\,h\,43\,m\,22\,s & 12 & C \\
23 & \emph{[in progress --- placeholder]} & --- & 31 & E \\
\bottomrule
\end{longtable}
}

\medskip
\noindent\footnotesize
\textbf{Status codes:}
\textbf{A} = OEIS A181671 published value.
\textbf{B} = prior self-computed, unpublished; cross-machine reproduced in this work (Table~\ref{tab:crossmachine}).
\textbf{C} = new result; both anchors cross-algorithm verified (Gourdon vs.\ Deleglise--Rivat).
\textbf{D} = new result; cross-algorithm verified \emph{and} independently reproduced via direct Miller--Rabin walk testing (Section~\ref{sec:verification}).
\textbf{E} = new result; computation ongoing at time of writing.
\normalsize

\subsection{Status of $a(23)$}
\label{sec:a23}

At the time of writing, $Q=10^{23}$ is being certified on the desktop
machine described in Section~\ref{sec:hardware}. All 31 grid-bracketing
checkpoints required for the safety net have been computed; the two
main anchors $\pcount(Q)$ and $\pcount(Q/2)$, each requiring independent
cross-verification via both the Gourdon and Deleglise--Rivat algorithms,
are in progress. \textbf{This placeholder will be replaced with the
certified value, its certificate, and full verification results once
the computation completes.} An independently published reference value
for $\pcount(10^{23})$, from the Wikipedia prime-counting-function
table~\cite{wikipedia-picount}, is available for cross-checking the
primary anchor once it lands: $\pcount(10^{23}) =
1{,}925{,}320{,}391{,}606{,}803{,}968{,}923$.

\section{Verification}
\label{sec:verification}

Every new term reported here (\(k=18\)--\(22\); $k=23$ pending) passes
the following independent checks, none of which are internal
self-consistency checks against the pipeline's own prior output:

\begin{enumerate}
  \item \textbf{Cross-algorithm anchor verification.} $\pcount(Q)$ and
  $\pcount(Q/2)$ for $k=20,21,22$ (and, redundantly, for the
  cross-machine reproduction of $k=18,19$) were each computed by two
  independently implemented algorithms within \texttt{primecount}
  (Gourdon and Deleglise--Rivat); all agreed exactly.
  \item \textbf{Independent external reference.} The Wikipedia
  prime-counting-function table~\cite{wikipedia-picount}, an
  independent source predating and unrelated to this project, lists
  $\pcount(10^{19})$ through $\pcount(10^{22})$; all four match this
  work's computed values digit-for-digit.
  \item \textbf{Miller--Rabin walk reproduction.} $a(21)$'s defining
  minimum occurs at an offset of 58 above $Q=10^{21}$, a dip of 3 below
  $f(Q)$. This was reproduced independently of \texttt{sieve128} by
  direct primality testing (via \texttt{sympy}) of every integer in
  $(Q, Q+58]$ and $(Q/2, (Q+58)/2]$: 0 primes were found in the former
  range and exactly 3 in the latter, confirming the dip via entirely
  different machinery. $a(20)$ and $a(22)$ have their minimum at offset
  0 (i.e.\ $a(k) = \pcount(Q)-\pcount(Q/2)$ exactly), so no walk
  verification is applicable to them.
  \item \textbf{Analytic agreement.} For $k=20,21,22$,
  $\li(Q)-\li(Q/2)$ agrees with the certified $f(Q)$ to within
  $6\times10^{-7}$ relative difference in every case.
  \item \textbf{Growth-ratio consistency.} $a(k)/a(k-1)$ for
  $k=19,\dots,22$ is $9.4699, 9.4966, 9.5208, 9.5427$ respectively ---
  a smooth, monotonically increasing sequence consistent with the
  expected $\sim\!9.5\times$ per-decade growth from prime density
  arguments, with no discontinuity at the transition into the
  cross-verified regime.
  \item \textbf{Analytic tail-bound formulas checked against source.}
  The Johnston (2022) and Dusart (2010) bounds as implemented were
  compared directly against the published statements in
  \cite{johnston2022,dusart2010}; both the constants and the stated
  domains of validity match exactly (Section~\ref{sec:background}).
  \item \textbf{Fresh, cache-bypassed single-pass recomputation.}
  $a(21)$ and $a(22)$ were originally computed across multiple
  interrupted sessions, so their first recorded wall-clock times
  reflected only a final resumed pass rather than a clean measurement.
  Both were independently recomputed from scratch, bypassing the
  persistent cache entirely, and matched the original certified values
  exactly, additionally producing publication-quality clean timings
  (Table~\ref{tab:results}).
  \item \textbf{Record-prime primality.} The largest Ramanujan prime
  below each certified bound (e.g.\ $99{,}999{,}999{,}999{,}999{,}999{,}989$
  for $k=20$) was independently confirmed prime via deterministic
  Miller--Rabin testing.
\end{enumerate}

\section{Performance analysis}
\label{sec:performance}

An initial analysis performed during this project, using timings
contaminated by interrupted and resumed runs under RAM pressure,
suggested per-decade cost growth as high as $11$--$26\times$, well above
the na\"ive $Q^{2/3} \approx 4.64\times$-per-decade asymptotic predicted
for the underlying algorithm. This estimate was an artifact of the
measurement conditions, not a property of the algorithm, and is
superseded here. Using the clean, single-pass timings in
Table~\ref{tab:results} (which for $k=18$--$22$ specifically bypass the
persistent cache and any partial-run contamination), a log-linear
regression of $\ln(\text{time})$ against $k$ gives:

\begin{itemize}
  \item $k=13$--$19$ (single-algorithm anchors): $\times3.51$ per decade of $Q$;
  \item $k=20$--$22$ (cross-verified anchors, roughly $3$--$4\times$
  slower per call than single-algorithm since Deleglise--Rivat is
  measurably slower than Gourdon at equal $Q$): $\times3.72$ per decade of $Q$.
\end{itemize}

Both regimes land in the same $3.5$--$3.7\times$-per-decade range despite
the methodology difference, corresponding to $\text{time} \propto
Q^{\,0.55\text{--}0.57}$ --- \emph{better} than the $Q^{2/3}$ asymptotic
usually quoted for this class of prime-counting algorithm. The number of
grid-bracketing checkpoints required grows far more slowly, consistent
with the $O(\log(W/D))$ prediction of Lemma~\ref{lem:bracket}, confirming
that per-call cost of each \texttt{primecount} evaluation, not the count
of evaluations needed, is what dominates the growth in total wall time.

\section{Limitations and the method's ceiling}
\label{sec:ceiling}

Two independent constraints bound how far this method can be pushed.

\emph{Mathematical ceiling.} The Johnston (2022) bound used for the
analytic tail is only valid for $y \leq 1.101\times10^{26}$. Beyond that,
this pipeline falls back to the weaker Dusart (2010) bounds, which remain
valid but require checking $\texttt{dusart\_ok}(m)$ can certify $f(y)
\geq m$ for all sufficiently large $y$ --- a condition that is expected
to fail once $Q$ approaches $1.101\times10^{26}$, placing the method's
absolute mathematical ceiling at approximately $a(26)$. Beyond that,
stronger published tail bounds would be required, not merely faster
hardware.

\emph{Practical ceiling.} Independent of the mathematical limit,
per-call RAM at $Q=10^{23}$ was measured at $\sim\!5.9$\,GB and is
increasing with an empirical exponent of roughly $0.5$--$0.7$ in $Q$
(Section~\ref{sec:implementation}), well above what fits comfortably
alongside an operating system and other workloads on consumer hardware
with 16--32\,GB of RAM. This was the direct cause of the mid-project
migration described in Section~\ref{sec:hardware}, and is expected to
bind again well before the mathematical ceiling at $a(26)$ is reached,
absent either a machine with substantially more RAM or algorithmic
improvements to \texttt{primecount}'s own memory footprint.

\section{Conclusion}

We have presented a certified, multiply-verified extension of OEIS
A181671 to $k=20,21,22$ (with $k=23$ in progress), computed entirely on
consumer-grade personal hardware. The approach combines an elementary
bracketing lemma with published, unconditional analytic tail bounds to
avoid sieving computationally infeasible ranges, backed by a
purpose-built 128-bit segmented sieve for the region that does require
direct primality information. Beyond the mathematical result, this work
required nontrivial systems engineering --- crash-safe persistent
caching, RAM-aware adaptive parallelism, and verified cross-machine
migration --- to complete reliably outside an institutional computing
environment, which we have documented here as it materially affects the
reproducibility of the reported timings and the practical ceiling on
how far the method can currently be pushed.

\section*{Acknowledgments}

The author thanks the maintainers of \texttt{primecount} and
\texttt{primesieve}, whose exact, high-performance implementations of
$\pcount(x)$ made this work possible.

\begin{thebibliography}{9}

\bibitem{ramanujan1919}
S.~Ramanujan, \emph{A proof of Bertrand's postulate}, Journal of the
Indian Mathematical Society \textbf{11} (1919), 181--182.

\bibitem{oeisA181671}
OEIS Foundation Inc., \emph{A181671: Number of Ramanujan primes} $\leq
10^n$, The On-Line Encyclopedia of Integer Sequences,
\url{https://oeis.org/A181671}.

\bibitem{dusart2010}
P.~Dusart, \emph{Estimates of some functions over primes without R.H.},
arXiv:1002.0442, 2010.

\bibitem{johnston2022}
D.~R.~Johnston, \emph{Improving bounds on prime counting functions by
partial verification of the Riemann hypothesis}, The Ramanujan Journal
(2022); arXiv:2109.02249.

\bibitem{primecount}
K.~Walisch et al., \emph{primecount: fast prime counting function
implementations}, \url{https://github.com/kimwalisch/primecount}.

\bibitem{primesieve}
K.~Walisch et al., \emph{primesieve: fast prime number generator},
\url{https://github.com/kimwalisch/primesieve}.

\bibitem{wikipedia-picount}
Wikipedia contributors, \emph{Prime-counting function}, Wikipedia, The
Free Encyclopedia. \url{https://en.wikipedia.org/wiki/Prime-counting_function}.

\end{thebibliography}

\end{document}
